\section{Methodology}
\label{sec:method}

\subsection{Problem Setup and Hypotheses}
For a model $m$ and a task $t$, we fine-tune on English and measure the test score $s_{m,t}(\ell)$ on each target language $\ell$.
The dependent variable is the \emph{transfer drop}
\begin{equation}
    \Delta_{m,t}(\ell) = s_{m,t}(\text{en}) - s_{m,t}(\ell),
    \label{eq:drop}
\end{equation}
in points, averaged over seeds.
Relative drop is a constant multiple of $\Delta$ within a model--task cell, so \rsq is identical.
The language is the unit of analysis within each of the six model--task cells.

We decompose the hypothesis into four pre-specified tests.
{\bf H1}: \rsq of $\Delta$ regressed on the fertility ratio lies in $[0.40, 0.60]$.
{\bf H2}: \rsq of fertility exceeds \rsq of typological distance.
{\bf H3}: languages with fertility ratio above 2.5 have $\Delta>30$, and languages below mostly do not.
{\bf H4}: the fertility effect survives controlling for pretraining data size.

\subsection{Models, Tasks and Languages}
\para{Models.}
We use \mbert (\texttt{bert-base-multilingual-cased}; WordPiece, 119k vocabulary; pretrained on Wikipedia in 104 languages) and \xlmr (\texttt{xlm-roberta-base}; SentencePiece unigram, 250k vocabulary; pretrained on CC-100 in 100 languages).
We choose these two because they are the standard zero-shot transfer baselines and they have different tokenizers, so a result is not a property of one vocabulary.

\para{Tasks.}
\Tabref{tab:setup} lists the three tasks.
We use three so that the result is not a property of one benchmark: \xnli~\cite{conneau2018xnli} trained on MNLI~\cite{williams2018mnli}, \wikiann NER~\cite{pan2017wikiann,rahimi2019massively}, and \sib topic classification~\cite{adelani2024sib200}.

\begin{table}[t]
    \centering
    \resizebox{0.9\textwidth}{!}{%
    \begin{tabular}{@{}lllllc@{}}
        \toprule
        Task & Dataset & English train & Test per language & Targets & Metric \\
        \midrule
        NLI & \xnli & 100k subset of MNLI & 5{,}010 & 14 & accuracy \\
        NER & \wikiann & 20{,}000 & 100 / 1{,}000 / 10{,}000 & 71 & span F1 \\
        Topic classification & \sib & 701 & 204 & 71 & accuracy \\
        \bottomrule
    \end{tabular}
    }
    \caption{Tasks. We train on English only and test zero-shot on every target language.}
    \label{tab:setup}
\end{table}


\para{Languages.}
We use all 72 languages (English plus 71 targets) that are in \mbert's pretraining list, in \xlmr's, in \flores and in \wikiann.
Restricting to languages that both models saw removes the unseen-language confound.

\para{Training.}
We use standard recipes with no per-language tuning: AdamW, a linear schedule with 10\% warmup, mixed precision, and model selection on English dev data only.
\xnli uses 3 seeds, NER 3 seeds and \sib 5 seeds, for 22 runs in total.
\Secref{sec:app_training} gives all hyperparameters.

\subsection{Predictors}
\para{Fertility ratio (primary).}
We count subword tokens per whitespace word on the language's 1{,}004 \flores sentences and divide by the same quantity for English (1.32 for \mbert, 1.40 for \xlmr).
Chinese, Japanese, Thai and Burmese are not whitespace-segmented, so we exclude them from the primary analysis.
This leaves 67 targets (12 for \xnli).

\para{Parallel token ratio (secondary).}
We divide the number of tokens on the \flores sentences by the number of tokens on their English translations.
This is defined for all 71 targets.

\para{Baselines.}
The main competitor is \uriel~\cite{littell2017uriel} syntactic distance from English, the typological distance used most in prior work.
We also report phonological, inventory, genetic and geographic distances, and all five jointly with adjusted \rsq.
Pretraining size is $\log_{10}$ bytes of the language's Wikipedia dump for \mbert and of its CC-100 file for \xlmr.
We also include a non-Latin script indicator.

\subsection{Statistical Analysis}
We report OLS \rsq with 95\% percentile bootstrap confidence intervals over languages (10{,}000 resamples), and a bootstrap interval for differences in \rsq.
For H4 we report partial correlations with OLS $t$-tests and the unique \rsq of each predictor in the model with fertility, log size and syntactic distance.
We apply Holm correction across the six cells for each family of tests.
For a pooled estimate we z-score the drop within each cell and cluster the bootstrap by language.
For H3 we compare the share of languages with $\Delta>30$ above and below 2.5$\times$ with Fisher's exact test, and we add a hinge at 2.5 to the linear fit as a check for nonlinearity.

\subsection{Two Designs That Hold the Language Fixed}
Across languages, fertility is correlated with size, typology and script.
We add two designs in which these cannot vary.

\para{Within-language contrast.}
For each language we relate the fertility difference between the two models (\mbert minus \xlmr) to the drop difference (\mbert minus \xlmr).
Typology and test data are identical on both sides.

\para{Over-segmentation intervention.}
We re-segment English test inputs so that no subword exceeds a cap of 6, 4, 3, 2 or 1 characters.
The text, the model and the labels are unchanged; only the granularity of the segmentation changes.
With no cap, our manual pipeline reproduces the standard tokenization exactly; we checked this on 300 examples per task and model and asserted it against the standard English score for every checkpoint.
This is the only experiment here that manipulates fertility directly.

\subsection{Deviations from the Plan}
We trained \xnli on a fixed 100k subset of MNLI instead of all 393k examples because of limited compute, so English accuracy is about 3 points below published full-data numbers.
We measure \mbert's pretraining size by Wikipedia dump bytes instead of article counts.
We added the script control in H4 after inspecting the fertility values but before running any regression.
\Secref{sec:app_deviations} lists all deviations.
